\documentclass[11pt]{article}

\usepackage[T1]{fontenc}
\usepackage{lmodern}
\usepackage{amsmath,amssymb,amsthm}
\usepackage[margin=1in]{geometry}
\usepackage{microtype}
\usepackage[hidelinks]{hyperref}
\hypersetup{
  pdftitle={Sharp pointwise bounds for Lebesgue functions},
  pdfauthor={Qiyuan Gu},
  pdfkeywords={Erdos Problem 1132, Lebesgue functions, polynomial
interpolation, almost everywhere, pointwise bounds, additive constants}
}

\newtheorem{theorem}{Theorem}
\newtheorem{lemma}[theorem]{Lemma}
\newtheorem{proposition}[theorem]{Proposition}
\newtheorem{taoestimates}[theorem]{Proposition}
\newtheorem{corollary}[theorem]{Corollary}
\theoremstyle{remark}
\newtheorem*{remark}{Remark}
\numberwithin{equation}{section}
\DeclareMathOperator{\supp}{supp}
\DeclareMathOperator{\sgn}{sgn}
\newcommand{\R}{\mathbb R}
\newcommand{\C}{\mathbb C}
\newcommand{\Z}{\mathbb Z}
\newcommand{\one}{\mathbf 1}

\title{Sharp pointwise bounds for Lebesgue functions}
\author{Qiyuan Gu\\
University of Chicago\\
\href{mailto:phoenix1203@uchicago.edu}{\texttt{phoenix1203@uchicago.edu}}}
\date{}

\begin{document}
\maketitle

\begin{abstract}
For each \(n\), let \(\lambda_n\) be the Lebesgue function for
polynomial interpolation at an arbitrary set of \(n\) distinct nodes
in \([-1,1]\). We prove that the points \(x\in(-1,1)\) for which
\(\lambda_n(x)>(2/\pi)\log n-C(x)\) infinitely often, with
some finite \(C(x)\), form a dense set. We also prove
\(\limsup_{n\to\infty}\lambda_n(x)/\log n\ge2/\pi\)
for almost every \(x\in(-1,1)\).
Conversely, we show that no absolute additive constant works for all arrays:
for every \(M>0\), there is an array for which each fixed
\(x\in(-1,1)\) eventually satisfies
\(\lambda_n(x)\le(2/\pi)\log n-M\).
For the same array, no single finite constant gives a dense set of
points with infinitely many lower-bound occurrences.
The positive results use derivative-jump energies and positive Cauchy
transforms. The counterexample combines an explicit logarithmic
integral estimate on the gaps of a compact set with real-rooted
polynomials having prescribed amplitudes.
\end{abstract}

\section{Introduction}

For each \(n\ge1\), let
\[
 -1\le x_{1,n}<\cdots<x_{n,n}\le1.
\]
Define the node polynomial, the cardinal polynomials, and the Lebesgue
function by
\begin{align*}
 P_n(x)&=\prod_{k=1}^n(x-x_{k,n}),\\
 \ell_{k,n}(x)&=\frac{P_n(x)}{(x-x_{k,n})P_n'(x_{k,n})},
 &\lambda_n(x)&=\sum_{k=1}^n|\ell_{k,n}(x)|.
\end{align*}
The cardinal polynomials are interpreted by continuity at the nodes.
The rows of the array need not be nested.

\begin{theorem}\label{thm:main}
For every triangular array of distinct nodes as above, the following
statements hold.
\begin{enumerate}
\item[(i)] The set of points \(x\in(-1,1)\) for which there is a finite
constant \(C(x)\) such that
\begin{equation}\label{eq:main-fixed}
 \lambda_n(x)>\frac2\pi\log n-C(x)
\end{equation}
for infinitely many \(n\) is dense in \((-1,1)\).
\item[(ii)] For Lebesgue almost every \(x\in(-1,1)\),
\begin{equation}\label{eq:main-ae}
 \limsup_{n\to\infty}\frac{\lambda_n(x)}{\log n}\ge\frac2\pi.
\end{equation}
\end{enumerate}
\end{theorem}

\begin{theorem}\label{thm:nonuniform}
For every \(M>0\), there is a triangular array of distinct nodes in
\((-1,1)\) such that, for every fixed \(x\in(-1,1)\),
\begin{equation}\label{eq:counter-eventual}
 \lambda_n(x)\le\frac2\pi\log n-M
 \qquad\text{for all sufficiently large }n.
\end{equation}
The starting index may depend on \(x\). The array can also be chosen
so that, for every finite \(C\), the set
\begin{equation}\label{eq:counter-good-set}
 G_C=\left\{x\in(-1,1):
 \lambda_n(x)>\frac2\pi\log n-C
 \text{ for infinitely many }n\right\}
\end{equation}
is not dense in \((-1,1)\).
\end{theorem}

Erd\H{o}s \cite[p.~68]{Erdos} asked whether every triangular array
has a point at which the sharp logarithmic lower bound holds
infinitely often with a bounded additive loss, and whether the
normalized lower bound \eqref{eq:main-ae} holds almost everywhere.
Theorem~\ref{thm:main}(i) answers the first question with a constant
that may depend on the point; part~(ii) answers the second.
As Tao observes in \cite[Remark~1.12]{Tao}, the original question does
not specify the allowed dependence of the additive constant.
Theorem~\ref{thm:nonuniform} distinguishes the uniform readings from
the point-dependent one.

For arbitrary triangular arrays, the three possible quantifier
patterns can now be distinguished, using \(G_C\) from
\eqref{eq:counter-good-set} for the array in question:
\begin{enumerate}
\item[(S1)] There is an absolute \(C\) such that \(G_C\ne\varnothing\)
for every array. This is false: use
Theorem~\ref{thm:nonuniform} with \(M>\max\{C,0\}\).
\item[(S2)] For every array, there is a finite \(C\), depending on the
array but not on the point, such that \(G_C\) is dense. This is false
by the second assertion of Theorem~\ref{thm:nonuniform}.
\item[(S3)] For every array, there are a finite \(C\) and a point
\(x\in G_C\). This is true by Theorem~\ref{thm:main}(i).
More precisely, \(\bigcup_{C\in\mathbb R}G_C\) is dense; the constant
may vary from point to point in this union.
\end{enumerate}

Theorem~\ref{thm:nonuniform} concerns arbitrary triangular arrays, as in
\cite[p.~68]{Erdos} and \cite[Corollary~1.11]{Tao}.
The problem page \cite{Bloom} instead uses the initial segments of a
single infinite node sequence. Our counterexample does not impose
this nesting, and the uniform-constant question in that setting
remains open.

\subsection{Background}

Faber \cite{Faber} proved the order \(\log n\) for the global Lebesgue
constant, Bernstein \cite{Bernstein} obtained the asymptotically sharp
leading constant \(2/\pi\), Erd\H{o}s and Tur\'an \cite{ErdosTuran}
improved the error to \(O(\log\log n)\), Erd\H{o}s \cite{Erdos1961}
to \(O(1)\), and V\'ertesi \cite[Corollary~1]{Vertesi1990} refined the
asymptotics through the optimal constant term to an error of
\(O((\log\log n/\log n)^2)\). See
\cite[Section~1.3]{Tao} and \cite{Brutman} for further historical discussion.

For triangular arrays, Erd\H{o}s \cite[p.~68]{Erdos} stated that the set
of points \(x\) satisfying
\[
 \limsup_{n\to\infty}\frac{\lambda_n(x)}{\log n}\ge\frac2\pi
\]
is everywhere dense, and asked whether this holds for almost every
\(x\). Erd\H{o}s and V\'ertesi \cite{ErdosVertesi} proved
almost-everywhere divergence of interpolants of a suitable continuous
function for every array. Their subsequent paper
\cite[Theorem~2.1]{ErdosVertesi1981} gives a quantitative lower bound
for the Lebesgue functions: for every \(\varepsilon>0\), there are
\(\eta(\varepsilon)>0\) and \(n_0(\varepsilon)\) such that, for every
array and \(n\ge n_0(\varepsilon)\), there is a measurable set
\(H_n\subset[-1,1]\) satisfying
\[
 |H_n|\le\varepsilon,\qquad
 \lambda_n(x)>\eta(\varepsilon)\log n
 \quad(x\in[-1,1]\setminus H_n).
\]
They note that Chebyshev nodes show this order to be optimal.
The exceptional sets may vary with \(n\); the estimate implies
\(\limsup_n\lambda_n(x)/\log n>0\) almost everywhere
(indeed, \(\left|\bigcup_N\bigcap_{n\ge N}H_n\right|\le\varepsilon\),
so outside a set of measure at most \(\varepsilon\), each point
leaves \(H_n\) infinitely often; now let \(\varepsilon\downarrow0\)).
Thus the logarithmic order was already known in 1981.
Theorem~\ref{thm:main}(ii) supplies the optimal coefficient \(2/\pi\),
which is the coefficient in Erd\H{o}s's question and in
\cite[Remark~1.12]{Tao}.
Erd\H{o}s and V\'ertesi also give the integral version
\cite[Corollary~2.2]{ErdosVertesi1981}:
\[
 \int_{S_n}\lambda_n(x)\,dx>
 (|S_n|-\varepsilon)\eta(\varepsilon)\log n
 \qquad(n\ge n_0(\varepsilon))
\]
for arbitrary measurable \(S_n\subset[-1,1]\).

Corollary~\ref{cor:positive-measure} gives, in every
sufficiently large row, a point in each fixed positive-measure set
with \(\lambda_n(x)>c\log n\), for any \(0<c<2/\pi\).
Erd\H{o}s--V\'ertesi's theorem instead controls all points outside a
small exceptional set, with coefficient \(\eta(\varepsilon)\).

Tao \cite[Theorem 1.10(i)]{Tao} proved that every fixed nondegenerate
interval \(I\subset[-1,1]\) satisfies
\[
 \sup_{x\in I}\lambda_n(x)\ge\frac2\pi\log n-O_I(1).
\]
His Corollary 1.11 and Remark 1.12 give, for every prescribed
\(\omega(n)\to\infty\), a dense, indeed comeager, set of points
\(x\) for which
\[
 \lambda_n(x)\ge\frac2\pi\log n-\omega(n)
\]
infinitely often. Theorem~\ref{thm:main}(i) replaces \(\omega(n)\) by a
finite constant depending on the point, on a dense set.
In footnote~6 to Remark~1.12 of \cite{Tao}, Tao expresses the tentative belief that
``this dependency on \(I\) in the error term cannot be entirely
eliminated.'' Corollary~\ref{cor:no-uniform-interval} establishes this
for the single array of Theorem~\ref{thm:nonuniform}.

\subsection{Organization}

Sections~\ref{sec:local}--\ref{sec:fixed} prove
Theorem~\ref{thm:main}(i). The local differentiation estimate turns
large derivative jumps into nearby high values of the Lebesgue
function. A weighted energy bound supplies enough such jumps, and a
second-moment argument gives recurrence on a positive-measure set.
Baire's theorem then gives density.

Sections~\ref{sec:cauchy}--\ref{sec:ae} prove
Theorem~\ref{thm:main}(ii) by a separate argument using positive
Cauchy transforms. This proof is independent of
Sections~\ref{sec:local}--\ref{sec:fixed} and of Tao's estimates.
Section~\ref{sec:counterexample} proves Theorem~\ref{thm:nonuniform}
independently of both lower-bound arguments. Its first three
subsections construct positive functions with large logarithmic
integral ratios; Lemma~\ref{lem:amplitude} converts them into
interpolation rows, which are assembled in the final subsection.
Section~\ref{sec:consequences} gives the positive-measure and
interval-constant corollaries.

Both main theorems and the two corollaries are formalized in Lean 4
with Mathlib. The estimates used from Tao in Proposition~\ref{prop:tao}
are proved directly in Lean. The file \path{Erdos1132/Theorems.lean}
collects the four statements. The \path{checks/} directory also gives them
with explicit node conditions and Lagrange products.
The source and build instructions are at
\url{https://github.com/FireflySentinel/erdos-1132}.

Throughout the paper, all logarithms are natural and \(|E|\) denotes
Lebesgue measure. We write \(c_0=2/\pi\), and suppress the row index
when discussing a single row. Constants in local estimates may depend
on the fixed intervals and test functions, but not on \(n\) or on
interpolation data bounded by one. The notation \(J\Subset U\) means
that \(J\) is a compact subset of the open set \(U\).

\section{A local differentiation formula}
\label{sec:local}


We use the following local form of Riesz's differentiation formula and its
near-equality consequence.

\begin{lemma}\label{lem:riesz}
Suppose that \(T\ge2\pi\) and \(F\) is holomorphic in a neighborhood of
the disk \(\{|z|\le T\}\), is real on its real diameter, and satisfies,
throughout this disk,
\[
 |F(u+iv)|\le e^{|v|}.
\]
Then
\begin{equation}\label{eq:riesz-derivative}
 |F'(0)|\le 1+O(T^{-1}).
\end{equation}
If \(F'(0)\ge 1-\delta\), where \(\delta\ge0\), then
\begin{equation}\label{eq:riesz-stability}
 F(\pi/2)\ge 1-\frac{\pi^2}{4}\delta-O(T^{-1}).
\end{equation}
The constants are absolute.

\end{lemma}

\begin{proof}
Choose \(R=m\pi\) with \(T/2\le R<T\), and integrate
\[
 \frac{F(z)}{z^2\cos z}
\]
around \(|z|=R\). On this circle, \(e^{|\Im z|}\le8|\cos z|\).
Indeed, for \(|v|\ge1\), use \(|\cos(u+iv)|\ge\sinh|v|\ge e^{|v|}/4\).
For \(|v|\le1\), put \(d=R-|u|=v^2/(R+|u|)\le1/2\); then
\(|\cos u|=|\cos d|\ge1/2\), while \(e^{|v|}\le e<3\).
The contour integral, divided by \(2\pi i\), has absolute value at most
\(8/R\). Taking residues at zero
and at \(t_j=(j+\tfrac12)\pi\) gives
\begin{equation}\label{eq:riesz-formula}
 F'(0)=\sum_{|t_j|<R}\frac{(-1)^jF(t_j)}{t_j^2}+O(T^{-1}).
\end{equation}
Applying the same formula to \(F(z)=\sin z\) gives
\[
 \sum_{|t_j|<R}\frac1{t_j^2}=1+O(T^{-1}).
\]
Together with \(|F(t_j)|\le1\), equation \eqref{eq:riesz-derivative} follows. Isolating the term \(t_0=\pi/2\)
in \eqref{eq:riesz-formula} gives
\[
 1-\delta\le \frac4{\pi^2}F(\pi/2)+1-\frac4{\pi^2}+O(T^{-1}),
\]
which proves \eqref{eq:riesz-stability}.
\end{proof}

We apply this formula to interpolants with bounded nodal data. We first
state the local estimates in the form used here, with the normalization:
\[
 S_n=\sum_k\frac1{|P_n'(x_{k,n})|},\qquad
 \alpha_n=\frac1n\log S_n,\qquad
 U_n(z)=\frac1n\log\frac1{|P_n(z)|}.
\]

\begin{taoestimates}[Tao's local estimates]\label{prop:tao}
Let \(I\Subset(-1,1)\) be a fixed nondegenerate closed interval, and
suppose that \(\sup_I\lambda_n\le n\) for all sufficiently large \(n\).
This is a specialization of the polynomial upper bound in
\cite[equation (4.1)]{Tao}. Define
\[
 \rho_n(x)=-\frac1{\pi^2}\int_{\mathbb R\setminus I}
                  \frac{U_n(y)-\alpha_n}{(y-x)^2}\,dy.
\]
For each fixed compact interval \(I'\Subset\operatorname{int}I\), the
following hold uniformly for all sufficiently large \(n\):
\begin{enumerate}
\item The densities are bounded above and below by positive constants
and have uniformly bounded Lipschitz norms:
\[
 \rho_n(x)\asymp_{I'}1,\qquad
 |\rho_n(x)-\rho_n(x_0)|\ll_{I'}|x-x_0|
 \qquad(x,x_0\in I').
\]
This is the form of \cite[Theorem 4.1(ii), equation (4.6)]{Tao} needed here.
\item For \(x\in I'\) and \(\eta\ge1/n\),
\begin{equation}\label{eq:potential}
 U_n(x+i\eta)=\alpha_n-\pi\eta\rho_n(x)
              +O_{I'}\!\left(\frac{\log n}{n}+\eta^3\right).
\end{equation}
Compare \cite[equation (4.7)]{Tao}.
\item For globally Lipschitz functions \(h_n\) supported in \(I'\), with
uniformly bounded supremum norms and Lipschitz constants,
\begin{equation}\label{eq:quadrature}
 \frac1n\sum_k h_n(x_{k,n})
 =\int h_n(x)\rho_n(x)\,dx+O(n^{-1/4}).
\end{equation}
\end{enumerate}
\end{taoestimates}

The Lean development proves this proposition directly in the stated form,
including the potential expansion
(\texttt{uniform\_\allowbreak{}complex\_\allowbreak{}potential\_\allowbreak{}expansion})
and density positivity
(\texttt{eventually\_\allowbreak{}exterior\_\allowbreak{}density\_\allowbreak{}positive}).
Thus the formal proof of Theorem~\ref{thm:main} does not depend on \cite{Tao}.

The density is defined from \(I\), as in \cite[equation (4.15)]{Tao}, and
is independent of the auxiliary interval \(I'\).
The quadrature estimate follows directly from the Poisson representation:
averaging Poisson kernels between heights \(h\) and \(2h\) gives an
approximate identity with Lipschitz error \(O(\sqrt h)\), while
\eqref{eq:potential} contributes \(O(\log n/(nh)+h^2)\).
Taking \(h=n^{-1/2}\) gives \eqref{eq:quadrature}.

\begin{lemma}\label{lem:local}
Let \(I\Subset(-1,1)\) be a fixed nondegenerate closed interval. Suppose that,
for all sufficiently large \(n\),
\begin{equation}\label{eq:local-upper}
 \lambda_n(x)\le \Lambda_n:=c_0\log n+C\qquad(x\in I),
\end{equation}
where \(C\) is fixed. This bound gives the hypothesis of the proposition
above; let \(\rho_n\) be its densities.
For every fixed interval \(J\Subset\operatorname{int}I\) and every
real interpolant
\[
 q(x)=\sum_{k=1}^n a_k\ell_{k,n}(x),\qquad |a_k|\le1,
\]
uniformly for \(x\in J\),
\begin{equation}\label{eq:derivative-bounds}
 |q'(x)|\le \pi n\rho_n(x)(\Lambda_n+o(1)),\qquad
 |q''(x)|\le C_Jn^2\log n.
\end{equation}
Moreover, for every fixed \(D\), if
\begin{equation}\label{eq:large-derivative}
 \frac{|q'(x)|}{\pi n\rho_n(x)}\ge c_0\log n-D,
\end{equation}
then there is \(y=x+O_J(n^{-1})\) such that
\begin{equation}\label{eq:nearby-peak}
 \lambda_n(y)\ge c_0\log n-D',
\end{equation}
where \(D'\) depends on \(I,J,C,D\), but not on \(n,x,q\). The point \(y\)
lies in any fixed compact interior interval containing \(J\) in its interior,
for all sufficiently large \(n\).
\end{lemma}

\begin{proof}
Fix \(I'\Subset\operatorname{int}I\) with \(J\Subset\operatorname{int}I'\).
For \(\Re z\in I'\) and \(\Im z=\eta\ge\log n/n\), the
interpolation formula and \eqref{eq:potential} give
\begin{equation}\label{eq:complex-interpolant}
 |q(z)|\le\frac{S_n|P_n(z)|}{\eta}
 \le\frac1\eta\exp\!\left(\pi n\eta\rho_n(\Re z)
                              +O(\log n+n\eta^3)\right).
\end{equation}
Fix \(x\in J\), and put
\[
 L=n^{-1/5},\qquad H=n^{-2/5},\qquad
 \omega=\pi n\rho_n(x)(1+n^{-1/10}).
\]
On the upper edge of \(\{|\Re z-x|\le L,\ 0\le\Im z\le H\}\),
we have \(|q(z)|\le\Lambda_n e^{\omega H}\): \(\omega H\) exceeds
\(\pi nH\rho_n(x)\) by \(\pi\rho_n(x)n^{1/2}\), which dominates the
\(O(n^{2/5})\) from the Lipschitz variation of \(\rho_n\) and the
remaining \(O(\log n)\) terms in \eqref{eq:complex-interpolant}.
On the lower edge \(|q|\le\Lambda_n\); a Chebyshev expansion on \(I\)
gives \(|q|\le\Lambda_n\exp(O_I(n))\) on the vertical sides.
For \(f(z)=\Lambda_n^{-1}q(z)e^{i\omega(z-x)}\), the horizontal
edges have \(|f|\le1\), and the vertical edges have \(|f|\le e^B\)
with \(B=O_I(n)\). Set
\[
 \mathcal B(z)=2B\left[e^{(z-x-L)/H}+e^{(-z+x-L)/H}\right].
\]
Its real part is nonnegative on the horizontal edges and at least
\(B\) on the vertical edges, since \(\cos(v/H)\ge1/2\) for \(0\le v\le H\).
The maximum modulus principle gives \(|f(z)e^{-\mathcal B(z)}|\le1\).
At distance at least \(L/2\) from the vertical sides, the resulting error
factor is \(\exp(O(n e^{-L/(2H)}))=1+O(n^{-10})\). Thus, also reflecting across
the real axis,
\begin{equation}\label{eq:local-growth}
 |q(x+u+iv)|\le\Lambda_n(1+n^{-10})e^{\omega|v|}
 \qquad(|u|\le L/2,\ |v|\le H).
\end{equation}
The rescaled function
\[
 F(t)=\frac{q(x+t/\omega)}{\Lambda_n(1+n^{-10})}
\]
satisfies Lemma~\ref{lem:riesz} on a disk of radius
\(T\asymp_J n^{3/5}\).
By \eqref{eq:riesz-derivative},
\[
 |q'(x)|=\omega\Lambda_n(1+n^{-10})|F'(0)|
 \le\pi n\rho_n(x)\Lambda_n(1+n^{-1/10})
                       (1+O(n^{-3/5})).
\]
Since \(\Lambda_n n^{-1/10}=o(1)\), this gives the first bound in
\eqref{eq:derivative-bounds}; Cauchy's estimate
on a disk of radius comparable to \(n^{-1}\) gives the second.
Under \eqref{eq:large-derivative}, change the sign of \(q\) if necessary.
For sufficiently large \(n\),
\[
 F'(0)\ge\frac{c_0\log n-D}
                 {\Lambda_n(1+n^{-1/10})(1+n^{-10})}
 \ge1-\frac{C_1}{\Lambda_n},\qquad C_1=1+\max\{C+D,0\}.
\]
Equation \eqref{eq:riesz-stability} then gives
\(|q(x+\pi/(2\omega))|\ge\Lambda_n-C_2\).
Since \(|q(y)|\le\lambda_n(y)\) on the real axis and \(\omega\asymp_J n\),
this proves \eqref{eq:nearby-peak} with the asserted displacement.
\end{proof}

\section{Derivative jumps and a Riesz energy}\label{sec:energy}

At each internal interpolation node, write
\[
 J_k=\lambda_n'(x_{k,n}+)-\lambda_n'(x_{k,n}-).
\]
At the first and last nodes, use the exterior polynomial pieces. Set
\[
 \nu_k=\frac1{S_n|P_n'(x_{k,n})|}.
\]
The absolute values of all cardinal polynomials other than \(\ell_{k,n}\) have
a cusp at \(x_{k,n}\), whereas \(\ell_{k,n}\) is positive near that node.
Using \(\ell_{j,n}(x)=P_n(x)/((x-x_{j,n})P_n'(x_{j,n}))\) and
\(P_n(x_{k,n})=0\) to differentiate at \(x_{k,n}\), we obtain
\begin{equation}\label{eq:jump}
 J_k=2\sum_{j\ne k}|\ell_{j,n}'(x_{k,n})|
     =2\sum_{j\ne k}\frac{\nu_j}{\nu_k|x_{k,n}-x_{j,n}|}.
\end{equation}

\begin{lemma}\label{lem:energy}
Under the hypotheses of Lemma~\ref{lem:local}, choose a nonzero nonnegative
smooth function \(f\) compactly supported in \(\operatorname{int}I\). Then
\begin{equation}\label{eq:jump-average}
 \sum_k\frac{f(x_{k,n})^2}{n\rho_n(x_{k,n})}
       \frac{J_k}{2\pi n\rho_n(x_{k,n})}
 \ge c_0\log n\int f(x)^2\,dx-O(1).
\end{equation}
Terms with \(f(x_{k,n})=0\) are taken to be zero; \(\rho_n\) is only needed on
a fixed interior neighborhood of the support of \(f\).

\end{lemma}

\begin{proof}
Put \(a_k=f(x_{k,n})/\rho_n(x_{k,n})\), with \(a_k=0\) outside the support of
\(f\), and define
\[
 \mu_n=\frac1n\sum_k a_k\delta_{x_{k,n}},\qquad
 K_\varepsilon(t)=\frac1{\sqrt{t^2+\varepsilon^2}}.
\]
The kernel \(K_\varepsilon\) is positive definite, since
\[
 K_\varepsilon(t)=\frac1{\sqrt\pi}
       \int_0^\infty s^{-1/2}e^{-s\varepsilon^2}e^{-st^2}\,ds
\]
is a positive mixture of Gaussian kernels. Write
\[
 \mathcal E_\varepsilon(\sigma,\tau)
 =\iint K_\varepsilon(x-y)\,d\sigma(x)\,d\tau(y)
\]
for its bilinear energy. Positivity applied to
\(\mu_n-f(x)\,dx\) gives
\begin{equation}\label{eq:energy-positivity}
 \mathcal E_\varepsilon(\mu_n,\mu_n)
 \ge2\mathcal E_\varepsilon(\mu_n,f\,dx)
       -\mathcal E_\varepsilon(f\,dx,f\,dx).
\end{equation}
Take \(\varepsilon=1/n\). Subtracting \(f(x)\) in the integral over
\(|t|\le1\), and using the Lipschitz bound for \(f\), gives the uniform
estimate
\begin{equation}\label{eq:kernel-asymptotic}
 (K_{1/n}*f)(x)=2\log n\,f(x)+O_f(1).
\end{equation}
The mass of \(\mu_n\) is bounded. Also, by \eqref{eq:quadrature},
\begin{equation}\label{eq:weighted-quadrature}
 \int f\,d\mu_n
 =\frac1n\sum_k\frac{f(x_{k,n})^2}{\rho_n(x_{k,n})}
 =\int f^2+O(n^{-1/10}).
\end{equation}
Again by \eqref{eq:kernel-asymptotic},
\[
 \mathcal E_{1/n}(f\,dx,f\,dx)=2\log n\int f^2+O(1).
\]
Together with \eqref{eq:energy-positivity}--\eqref{eq:weighted-quadrature}, this
yields
\begin{equation}\label{eq:energy-lower}
 \mathcal E_{1/n}(\mu_n,\mu_n)
 \ge2\log n\int f^2-O(1).
\end{equation}
The diagonal part of this energy is \(n^{-1}\sum_k a_k^2=O(1)\). Since
\(K_{1/n}(t)\le1/|t|\) for \(t\ne0\),
\begin{equation}\label{eq:off-diagonal}
 \frac1{n^2}\sum_{k\ne j}\frac{a_ka_j}{|x_{k,n}-x_{j,n}|}
 \ge2\log n\int f^2-O(1).
\end{equation}
Finally,
\[
 a_k^2\frac{\nu_j}{\nu_k}+a_j^2\frac{\nu_k}{\nu_j}
 \ge2a_ka_j.
\]
By \eqref{eq:jump}, the left side of \eqref{eq:jump-average} is
\[
 \frac1{\pi n^2}\sum_{k\ne j}
       \frac{a_k^2\nu_j}{\nu_k|x_{k,n}-x_{j,n}|}.
\]
Pairing the terms indexed by \((k,j)\) and \((j,k)\), and then applying
\eqref{eq:off-diagonal}, bounds this expression below by
\[
 \frac1{\pi n^2}\sum_{k\ne j}
       \frac{a_ka_j}{|x_{k,n}-x_{j,n}|}
 \ge\frac2\pi\log n\int f^2-O(1).
\]
This proves \eqref{eq:jump-average}.
\end{proof}

\section{A dense set of points with bounded additive loss}\label{sec:fixed}

\begin{proposition}\label{prop:recurrence}
Assume the local upper bound \eqref{eq:local-upper} on a fixed
nondegenerate interval \(I\Subset(-1,1)\). There are a constant \(D\)
and a measurable set \(G\subset\operatorname{int}I\), with \(|G|>0\),
such that every \(x\in G\) satisfies
\[
 \lambda_n(x)\ge c_0\log n-D
\]
for infinitely many \(n\).
\end{proposition}

\begin{proof}
Choose compact intervals \(J_0\Subset J_1\Subset\operatorname{int}I\), and
choose \(f\) as in Lemma~\ref{lem:energy} with \(\supp f\subset J_0\).
All constants below are independent of \(n\): after \(I,C,J_0,J_1,f\) are
fixed, we choose \(D\), then \(\kappa,D_1,d,\kappa_1,d_1\), and finally
\(b,D_2\). Let
\[
 w_k=\frac{f(x_{k,n})^2}{n\rho_n(x_{k,n})},\qquad
 T_k=\frac{J_k}{2\pi n\rho_n(x_{k,n})},\qquad W=\int f^2>0.
\]
By \eqref{eq:quadrature},
\begin{equation}\label{eq:weights}
 \sum_kw_k=W+O(n^{-1/10}),\qquad 0\le w_k\le C_f/n.
\end{equation}
On each interval between successive nodes, \(\lambda_n\) equals an interpolant
with nodal data in \(\{-1,1\}\). Both one-sided derivatives at a node
therefore satisfy \eqref{eq:derivative-bounds}, so
\begin{equation}\label{eq:jump-upper}
 0\le T_k\le\Lambda_n+o(1)
\end{equation}
uniformly on the support of \(f\). Lemma~\ref{lem:energy}, \eqref{eq:weights},
and \eqref{eq:jump-upper} show that
\begin{equation}\label{eq:weighted-deficit}
 \sum_kw_k(\Lambda_n+1-T_k)=O(1),
\end{equation}
with nonnegative summands for all sufficiently large \(n\). On a node with
\(T_k<c_0\log n-D\), the factor in \eqref{eq:weighted-deficit} is greater
than \(C+1+D\). Hence, by choosing \(D\) sufficiently large, such nodes
carry at most \(W/2\) of the weights. Since \eqref{eq:weights} has total
weight \(W+o(1)\), for all
sufficiently large \(n\) the nodes in the support of \(f\) satisfying
\begin{equation}\label{eq:large-jump}
 T_k\ge c_0\log n-D
\end{equation}
carry at least \(W/4\) of the weights. In particular, there are at least
\(\kappa n\) such nodes, with \(\kappa>0\) independent of \(n\). Discarding the first and last nodes does not change this conclusion
after reducing \(\kappa\).

At a node satisfying \eqref{eq:large-jump}, at least one adjacent polynomial
piece \(q\) has
\[
 |q'(x_{k,n})|\ge\pi n\rho_n(x_{k,n})(c_0\log n-D).
\]
Lemma~\ref{lem:local} gives a point
\begin{equation}\label{eq:high-points}
 y_k=x_{k,n}+O(n^{-1}),\qquad
\lambda_n(y_k)\ge c_0\log n-D_1.
\end{equation}
The same polynomial piece \(q\) has value 1 at both endpoints of the
corresponding gap. Rolle's theorem and
\eqref{eq:derivative-bounds} imply that this gap has length at least \(d/n\)
for some fixed \(d>0\) when the gap is contained in \(J_1\). If it is not
contained in \(J_1\), its length is bounded below by
\(\operatorname{dist}(J_0,\mathbb R\setminus J_1)>0\).

An interval of length less than \(d/n\) contains at most two such nodes, so
every interval of length \(O(n^{-1})\) contains only a bounded number of the
nodes in \eqref{eq:large-jump}. By \eqref{eq:high-points}, the same holds,
with multiplicities, for the points \(y_k\). Hence one can select at least
\(\kappa_1n\) of them, separated by \(d_1/n\), for fixed
\(\kappa_1,d_1>0\).

Equation \eqref{eq:derivative-bounds}, applied to every polynomial piece, also
shows that \(\lambda_n\) is \(C_In\log n\)-Lipschitz on a fixed compact
interior interval containing these points. Choose a fixed sufficiently small
\(b>0\), and let \(A_n\) be the union of the intervals of radius
\[
 r_n=\frac{b}{n\log n}
\]
centered at the selected points. For all sufficiently large \(n\), these
intervals are disjoint, lie in a fixed compact interior interval, and satisfy
\begin{equation}\label{eq:high-sets}
 \lambda_n(x)\ge c_0\log n-D_2\quad(x\in A_n),\qquad
 \frac{c}{\log n}\le |A_n|\le\frac{C_2}{\log n}.
\end{equation}
The separation of the centers in the \(m\)-th row implies, for every interval
\(B\),
\[
 |A_m\cap B|\le C_3\left(\frac{|B|}{\log m}
                              +\frac1{m\log m}\right).
\]
Summing this inequality over the at most \(n\) components of \(A_n\)
gives, for \(m>n\),
\begin{equation}\label{eq:intersection}
 |A_n\cap A_m|
 \le C_3\left(\frac{|A_n|}{\log m}
                              +\frac{n}{m\log m}\right).
\end{equation}
Put \(B_j=A_{2^j}\). Equations \eqref{eq:high-sets} and
\eqref{eq:intersection} give
\begin{equation}\label{eq:dyadic-intersection}
 |B_j|\asymp j^{-1},\qquad
 |B_j\cap B_k|\le C_4\left(\frac1{jk}+\frac{2^{j-k}}k\right)
 \quad(j<k).
\end{equation}
For every fixed \(j_0\), the sum of the measures of \(B_j\), \(j_0\le j\le
N\), grows at least as a positive constant times \(\log N\). By
\eqref{eq:dyadic-intersection}, the integral of the square of the sum of their
indicators is \(O((\log N)^2)\). Cauchy--Schwarz consequently gives a positive
lower bound, independent of \(j_0\), for
\[
 \left|\bigcup_{j\ge j_0}B_j\right|.
\]
Continuity of measure along these decreasing tail unions shows that
\begin{equation}\label{eq:recurrence}
 |\limsup_{j\to\infty}B_j|>0.
\end{equation}
Together with \eqref{eq:high-sets}, this proves the proposition.
\end{proof}

\begin{proof}[Proof of Theorem~\ref{thm:main}(i)]
Fix a nondegenerate compact interval \(K_0\subset(-1,1)\), and suppose
that it contains no point satisfying the assertion. Then
\begin{equation}\label{eq:additive-failure}
 \lambda_n(x)-c_0\log n\longrightarrow-\infty
 \qquad\text{for every }x\in K_0.
\end{equation}
The closed sets
\[
 F_M=\{x\in K_0:\lambda_n(x)\le c_0\log n+M
                   \text{ for every }n\ge2\},
 \qquad M=1,2,\ldots,
\]
cover \(K_0\), by \eqref{eq:additive-failure}. Baire's theorem supplies
a nondegenerate closed subinterval \(I\subset K_0\) on which
\eqref{eq:local-upper} holds with a fixed \(C\).
Proposition~\ref{prop:recurrence} contradicts
\eqref{eq:additive-failure}. Thus every such \(K_0\) contains a good
point, proving density. Increasing the final additive constant by one
gives the strict inequality in the statement.
\end{proof}

\section{Positive Cauchy transforms}\label{sec:cauchy}

Suppress the row index and put
\[
 S=\sum_k\frac1{|P'(x_k)|},\qquad
 \nu_k=\frac1{S|P'(x_k)|},\qquad
 \epsilon_k=\operatorname{sgn}P'(x_k).
\]
Define
\begin{equation}\label{eq:transforms}
 m(z)=\sum_k\frac{\nu_k}{z-x_k},\qquad
 g(z)=\sum_k\frac{\epsilon_k\nu_k}{z-x_k}
     =\frac1{SP(z)}.
\end{equation}
The last identity is the partial fraction formula for \(1/P\), whose expansion
at infinity gives
\begin{equation}\label{eq:moments}
 \sum_k\epsilon_k\nu_kx_k^j=0\qquad(0\le j\le n-2).
\end{equation}
For \(n\ge2\), both \(m+g\) and \(m-g\) are Cauchy transforms of positive
probability measures: their masses at \(x_k\) are \((1\pm\epsilon_k)\nu_k\).

For real \(x\) away from the nodes, write
\[
 H(x)=\sum_k\frac{\nu_k}{|x-x_k|},\qquad
 \gamma_s(x)=-\Im m(x+is),\qquad
 p_s(t)=\frac{s}{\pi(t^2+s^2)}.
\]
Then
\begin{equation}\label{eq:lebesgue-transform}
 \lambda_n(x)=\frac{H(x)}{|g(x)|},\qquad
 H(x)=\frac2\pi\int_0^\infty\gamma_s(x)\,\frac{ds}{s},
\end{equation}
and the Poisson semigroup gives
\begin{equation}\label{eq:semigroup}
 \gamma_s=p_{s-h}*\gamma_h\qquad(s>h).
\end{equation}

\begin{lemma}\label{lem:cancellation}
For every row with \(n\ge2\), every \(0<h\le1\), and every
\(x\in[-1,1]\),
\begin{equation}\label{eq:cancellation-bound}
 \frac{|g(x+ih)|}{\gamma_h(x)}
 \le\frac{4+h^2}{h^2\sinh((n-1)h/2)}.
\end{equation}
Consequently, if \(0<a<1\) and \(h=n^{-a}\), then
\begin{equation}\label{eq:cancellation}
 \sup_{x\in[-1,1]}\frac{|g(x+ih)|}{\gamma_h(x)}\longrightarrow0.
\end{equation}
\end{lemma}

\begin{proof}
Set \(z=x+ih\) and \(p=T_{n-1}\), the Chebyshev polynomial of the
first kind. The polynomial \((p(t)-p(z))/(t-z)\) has degree at most
\(n-2\), so \eqref{eq:moments} gives
\[
 p(z)g(z)=\sum_k\frac{\epsilon_k\nu_kp(x_k)}{z-x_k},
 \qquad |p(z)g(z)|\le\frac1h,
\]
using \(|p(x_k)|\le1\) and \(\sum_k\nu_k=1\).
Write \(z=\cos(u+iv)\), with \(u,v\in\R\). Then
\(h\le\sinh|v|\), so \(|v|\ge\operatorname{arsinh}h\ge h/2\).
The identity \(T_{n-1}(\cos w)=\cos((n-1)w)\) yields
\[
 |p(z)|^2=\cos^2((n-1)u)+\sinh^2((n-1)v)
 \ge\sinh^2((n-1)h/2).
\]
Finally, since the nodes and \(x\) lie in \([-1,1]\),
\[
 \frac{h}{4+h^2}\le
 \gamma_h(x)=\sum_k\frac{\nu_kh}{(x-x_k)^2+h^2}.
\]
Combining these bounds proves \eqref{eq:cancellation-bound}.
For \(h=n^{-a}\), its right-hand side tends to zero: the hyperbolic
sine grows exponentially in \(n^{1-a}\), while \(h^{-2}\) grows
polynomially. This proves \eqref{eq:cancellation}.
\end{proof}

\begin{lemma}\label{lem:boundary}
Let \(F\) be the Cauchy transform of a positive probability measure with
finite support. For every bounded real interval \(S_0\), every \(x\in\R\), and
every \(h>0\),
\begin{equation}\label{eq:harmonic-measure}
 \int_{\mathbb R}p_h(x-y)\,
            \mathbf1_{\{F(y)\in S_0\}}\,dy
 =\frac1\pi\int_{S_0}
       \frac{-\Im F(x+ih)}{(t-\Re F(x+ih))^2+(\Im F(x+ih))^2}\,dt.
\end{equation}
\end{lemma}

\begin{proof}
The function \(F\) maps the upper half-plane into the lower half-plane.
Compose harmonic measure of \(S_0\) there with \(F\). The
resulting bounded harmonic function on the upper half-plane has the indicated
boundary values almost everywhere. The exceptional boundary points are poles
of \(F\) and preimages of the endpoints of \(S_0\), a finite set. Its Poisson
representation is \eqref{eq:harmonic-measure}.
\end{proof}

\section{The almost-everywhere bound}\label{sec:ae}

\begin{proof}[Proof of Theorem~\ref{thm:main}(ii)]
The failure set is the union, over rational \(0<c<c_0\) and integers
\(N\ge2\), of the measurable sets
\[
 \{y\in(-1,1):\lambda_n(y)\le c\log n\text{ for every }n\ge N\}.
\]
If the assertion fails, one of these sets \(E\) has positive measure.
For its fixed \(c,N\), we therefore have
\begin{equation}\label{eq:common-low-set}
 \lambda_n(y)\le c\log n\qquad(y\in E,\ n\ge N).
\end{equation}
Choose constants
\[
 K>1,\qquad 0<b<a<1,\qquad \frac{2b}{\pi cK}>1,
\]
which is possible because \(c<2/\pi\). Set
\[
 h=n^{-a},\qquad\eta=n^{-b},\qquad
 H^\eta(x)=\frac2\pi\int_\eta^1\gamma_s(x)\,\frac{ds}{s}.
\]
By \eqref{eq:lebesgue-transform}, \(H^\eta\le H\). For \(s>0\)
and real \(x,y,t\), the triangle inequality gives
\[
 |y+is-t|\le |x+is-t|+|x-y|
 \le\left(1+\frac{|x-y|}{s}\right)|x+is-t|.
\]
Taking reciprocal squares and summing against the positive masses yields
\(\gamma_s(x)\le(1+|x-y|/s)^2\gamma_s(y)\). Integrating over
\(\eta\le s\le1\), we obtain
\begin{equation}\label{eq:potential-comparison}
 H^\eta(x)\le\left(1+\frac{|x-y|}{\eta}\right)^2H^\eta(y).
\end{equation}

Fix \(x\in E\), and write
\[
 u=\Re m(x+ih),\qquad\gamma=\gamma_h(x),\qquad
 S_x=[u-K\gamma,u+K\gamma].
\]
By \eqref{eq:cancellation} and \eqref{eq:harmonic-measure}, each of the two events
\[
 m(y)+g(y)\in S_x,\qquad m(y)-g(y)\in S_x
\]
has probability
\[
 p_K+o(1),\qquad p_K=\frac2\pi\arctan K>\frac12,
\]
under the Poisson probability \(p_h(x-y)\,dy\). The error is uniform in \(x\in
E\): for \(F=m\pm g\), \eqref{eq:cancellation} gives
\(\Re F(x+ih)=u+o(\gamma)\) and
\(-\Im F(x+ih)=\gamma(1+o(1))\). Substituting \(t=u+\gamma s\)
in \eqref{eq:harmonic-measure} therefore gives
\(\pi^{-1}\int_{-K}^K(1+s^2)^{-1}\,ds+o(1)=p_K+o(1)\).
The two events have intersection probability at least
\begin{equation}\label{eq:event-intersection}
 q_K+o(1),\qquad q_K=2p_K-1>0.
\end{equation}
Choose \(R\) so that the Poisson mass of \(|y-x|>Rh\) is less than \(q_K/8\),
and let
\[
 E_n=\{x\in E:(p_h*\mathbf1_{\mathbb R\setminus E})(x)<q_K/8\}.
\]
The approximate identity theorem gives \(|E\setminus E_n|\to0\). For all
sufficiently large \(n\), \eqref{eq:event-intersection} is greater than
\(q_K/2\). For \(x\in E_n\), the two excluded sets
\(\R\setminus E\) and \(\{|y-x|>Rh\}\) have total Poisson mass
less than \(q_K/4\). Hence there exists \(y\in E\) with
\(|y-x|\le Rh\) and both events true. We may take \(y\) away from
the nodes. Since \(m(y)\pm g(y)\) both lie in an interval of length
\(2K\gamma_h(x)\),
\[
 |g(y)|\le K\gamma_h(x).
\]
Equations \eqref{eq:lebesgue-transform} and \eqref{eq:common-low-set} give
\[
 H^\eta(y)\le H(y)=\lambda_n(y)|g(y)|
 \le cK(\log n)\gamma_h(x).
\]
Together with \eqref{eq:potential-comparison}, this yields
\begin{equation}\label{eq:local-comparison}
 H^\eta(x)\le(1+Rh/\eta)^2H^\eta(y)
 \le cK(1+Rh/\eta)^2(\log n)\gamma_h(x)
 \qquad(x\in E_n).
\end{equation}

Let \(T_n\) be convolution with the symmetric probability kernel
\[
 k_n(t)=\frac1{b\log n}\int_\eta^1p_{s-h}(t)\,\frac{ds}{s}.
\]
Its integral is one, since \(\int_\eta^1 ds/s=b\log n\).
Equation \eqref{eq:semigroup} then shows that
\[
 H^\eta=\frac{2b}{\pi}(\log n)T_n\gamma_h.
\]
Thus \eqref{eq:local-comparison} becomes
\begin{equation}\label{eq:kernel-ratio}
 \frac{T_n\gamma_h(x)}{\gamma_h(x)}\le\frac1{A_n}
 \quad(x\in E_n),\qquad
 A_n=\frac{2b}{\pi cK}(1+Rh/\eta)^{-2}\longrightarrow A>1.
\end{equation}
Integrate \eqref{eq:kernel-ratio} over \(E_n\), discard the contribution from
\(y\notin E_n\), and symmetrize. Positivity of \(\gamma_h\), symmetry of
\(k_n\), and \((r+r^{-1})/2\ge1\) give
\begin{equation}\label{eq:symmetric-energy}
 \frac{|E_n|}{A_n}
 \ge\iint_{E_n\times E_n}k_n(x-y)
                  \frac{\gamma_h(y)}{\gamma_h(x)}\,dy\,dx
 \ge\iint_{E_n\times E_n}k_n(x-y)\,dy\,dx.
\end{equation}
The kernels \(k_n\) form an approximate identity. In fact, for every fixed
\(\delta>0\),
\[
 \int_{|t|>\delta}k_n(t)\,dt
 \le\frac{C}{b\log n}\int_\eta^1\min(1,s/\delta)\,\frac{ds}{s}
 \longrightarrow0.
\]
It follows that
\[
 \iint_{E\times E}k_n(x-y)\,dy\,dx\longrightarrow|E|.
\]
Replacing \(E\) by \(E_n\) changes this integral by at most \(2|E\setminus
E_n|\). Taking limits in \eqref{eq:symmetric-energy} therefore yields
\[
 \frac{|E|}{A}\ge|E|,
\]
contrary to \(|E|>0\) and \(A>1\). This proves the assertion.
\end{proof}



\section{Failure of uniform additive constants}
\label{sec:counterexample}

We prove Theorem~\ref{thm:nonuniform}. The constants are chosen
for a clean verification and are not optimized. Smooth on \(I=[-1,1]\)
will mean extendible to a smooth function on a neighborhood of \(I\).

\subsection{The compact set}\label{subsec:counter-set}

For \(I=[-1,1]\), write
\[
 Lf(x)=\int_I\frac{f(x)-f(y)}{|x-y|}\,dy.
\]
In particular,
\begin{equation}\label{eq:counter-u3}
\|Lf\|_\infty\le2\|f'\|_\infty\qquad(f\in C^1(I)).
\end{equation}

Fix \(A>1\), and put
\[
 \rho=e^{-1024(A+40)},\qquad
 \delta_k=\rho\,2^{-k},\qquad
 m_k=\left\lceil e^{1/\delta_k}\right\rceil .
\]
Starting with \(I\), replace every retained interval of length \(\ell_{k-1}\) at stage \(k\) by \(m_k\) equally long closed intervals, retaining both endpoints of the parent. The children and intervening open gaps have respective lengths
\[
 \ell_k=\frac{(1-\delta_k)\ell_{k-1}}{m_k},\qquad
 b_k=\frac{\delta_k\ell_{k-1}}{m_k-1},\qquad \ell_0=2.
\]
Let \(F\) be the intersection of the retained sets and \(U=I\setminus F\). Thus \(F\) is compact, contains both endpoints, and has empty interior, since \(\ell_k\to0\). The set \(U\) is a dense open union of disjoint gaps, and
\begin{equation}\label{eq:counter-u4}
|U|\le2\sum_k\delta_k=2\rho.
\end{equation}

We need the following two properties:
\begin{equation}\label{eq:counter-u5}
|F\cap(x-r,x+r)|\ge\kappa r
 \quad(x\in F,\ 0<r\le1),\qquad \kappa=1/32,
\end{equation}
and
\begin{equation}\label{eq:counter-u6}
\int_U\frac{dy}{|x-y|}=\infty\qquad(x\in F).
\end{equation}

For \eqref{eq:counter-u5}, choose the first \(k\) for which \(\ell_k\le r/16\), and let \(J\) be the stage \(k-1\) parent containing \(x\), with \(|J|>r/16\). At least half the length of any retained interval belongs to \(F\), since every subsequent product of retained proportions is at least \(1-\rho>1/2\). If \(|J|\le r\), then \(J\subset[x-r,x+r]\), proving \eqref{eq:counter-u5}. Otherwise \(J\cap[x-r,x+r]\) has length \(s\ge r\). Write \(a=\ell_k\), \(b=b_k\), \(\delta=\delta_k\). Then
\[
 b\le a,\qquad b/(a+b)\le2\delta.
\]
The gaps intersecting an interval of length \(s\) have total length at most \(2\delta s+2b\). The two partially intersected children have total length at most \(2a\). The full children inside this interval therefore have total length at least
\[
 (1-2\delta)s-2(a+b)\ge r/2.
\]
Their intersections with \(F\) have total length at least \(r/4\), as required.

For \eqref{eq:counter-u6}, in the stage \(k-1\) parent containing \(x\), one side of its child contains at least \((m_k-1)/2\) gaps. Use the first \(q=\lfloor m_k/8\rfloor\) on that side. The far endpoint of the \(j\)-th such gap is at distance at most \(4j\ell_{k-1}/m_k\) from \(x\), and its length is at least \(\delta_k\ell_{k-1}/m_k\). These gaps contribute at least
\[
 \frac{\delta_k}{4}\sum_{j=1}^{q}\frac1j
 \ge\frac{\delta_k}{4}(\log m_k-\log16)\ge\frac18.
\]
The gaps used at different stages are disjoint. Summation proves \eqref{eq:counter-u6}.

\subsection{An explicit logarithmic integral estimate}\label{subsec:counter-integral}

On each gap \((b,c)\), define
\[
 r(x)=\frac{(x-b)(c-x)}{c-b},
\]
and put \(r=0\) on \(F\). If \(D(x)=\operatorname{dist}(x,F)\), then
\begin{equation}\label{eq:counter-u7}
r(x)\le D(x)\le2r(x)\qquad(x\in U).
\end{equation}
The function \(r\) is globally \(1\)-Lipschitz: its derivative has absolute value at most one on each gap, and points in distinct gaps can be compared through a point of \(F\) between them.

Put
\begin{equation}\label{eq:counter-u8}
B=e^8,\qquad \alpha=1/256,\qquad
 p(t)=[\log(B/t)]^{-\alpha},\qquad
 u(x)=p(r(x))\ (x\in U),\quad u=0\ (x\in F).
\end{equation}
Then \(u\) is continuous on \(I\), bounded, positive and smooth in each gap, and vanishes on \(F\). In particular, \(Lu(x)\) is an absolutely convergent integral for every \(x\in U\).

We prove the quantitative bound
\begin{equation}\label{eq:counter-u9}
\frac{Lu(x)}{u(x)}
 \ge \frac1{512}\log\frac1{4r(x)}-32
 \qquad\text{if }\log(B/r(x))\ge16.
\end{equation}
All estimates below apply uniformly up to the endpoints of \(I\).

Fix such an \(x\), and write \(s=r(x)\), \(T=\log(B/s)\), \(a_0=4s\).
For \(t\in[0,2]\), let
\[
 M_I(t)=|I\cap(x-t,x+t)|,\qquad
 M_U(t)=|U\cap(x-t,x+t)|.
\]
Choose \(z\in F\) with \(|x-z|=D(x)\le2s\). If \(4s\le t\le2\), the ball of radius \(t/2\) about \(z\), intersected with \(I\), is contained in the ball of radius \(t\) about \(x\). By \eqref{eq:counter-u5}, the latter contains at least \(\kappa t/2\) of \(F\). Since \(M_I(t)\le2t\),
\begin{equation}\label{eq:counter-u10}
M_U(t)\le(1-\eta)M_I(t),\qquad \eta=\kappa/4=1/128
 \quad(4s\le t\le2).
\end{equation}

First consider \(|y-x|<4s\). For \(|y-x|<s/2\), Lipschitz continuity of \(r\) gives
\(r(y)\in[s/2,3s/2]\). On this range \(p'(r(y))\le4p(s)/s\), by the derivative formula for \(p\) and \(T\ge16\). The integral of the absolute difference quotient over this part is thus at most \(4p(s)\). On \(s/2\le|y-x|<4s\), we have \(r(y)\le5s\) and
\[
 p(5s)/p(s)=\left(\frac{T}{T-\log5}\right)^\alpha\le2.
\]
The absolute difference is at most \(3p(s)\), so this part contributes at most \(6\log8\,p(s)<18p(s)\). In particular, the whole near contribution to \(Lu(x)\) is at least \(-24p(s)\).

For the far contribution, set \(l_-=1+x\), \(l_+=1-x\), and
\[
 H=\sum_{\nu\in\{-,+\}}\max\{\log(l_\nu/(4s)),0\}.
\]
The positive part is exactly \(p(s)H\). For the negative part,
\(r(y)\le s+|x-y|\le(5/4)|x-y|\), so it is bounded by
\[
 \int_{a_0}^2 f(t)\,dM_U(t),\qquad
 f(t)=\frac{p(5t/4)}t.
\]
The function \(f\) is decreasing. Stieltjes integration by parts, dropping the nonpositive lower-end boundary term and using \eqref{eq:counter-u10}, gives
\begin{equation}\label{eq:counter-u11}
\int_{a_0}^2 f(t)\,dM_U(t)
 \le(1-\eta)\left[f(a_0)M_I(a_0)+
                     \int_{a_0}^2 f(t)\,dM_I(t)\right].
\end{equation}
The first term on the right, including its factor, is at most \(4p(s)\).

Since \(dM_I(t)=1_{\{t<l_-\}}dt+1_{\{t<l_+\}}dt\), the remaining term splits into at most two integrals. For a side with \(l_\nu>4s\), put
\[
 h=\log(l_\nu/(4s)),\qquad T_0=T-\log5.
\]
Then \(0\le h<T_0\), since \(T_0-h=\log(4B/(5l_\nu))>0\), and
\begin{equation}\label{eq:counter-u12}
\int_{4s}^{l_\nu}\frac{p(5t/4)}t\,dt
 =\int_0^h(T_0-z)^{-\alpha}\,dz
 \le\frac{T_0^{-\alpha}}{1-\alpha}\,h.
\end{equation}
The last inequality follows from
\((1-z)^{1-\alpha}\ge1-z\) for \(0\le z\le1\).

Finally, concavity of the power \(\alpha\), \(\log5<2\), and \(T\ge16\) give
\begin{equation}\label{eq:counter-u13}
\begin{split}
 \frac{1-\eta}{1-\alpha}\frac{T_0^{-\alpha}}{p(s)}
 &=\frac{254}{255}\left(\frac{T}{T-\log5}\right)^{1/256}\\
 &\le\frac{254}{255}(8/7)^{1/256}
 \le\frac{254}{255}\left(1+\frac1{1792}\right)
 \le1-\frac1{512}.
\end{split}
\end{equation}
Thus the negative far contribution is at most
\(4p(s)+(1-1/512)p(s)H\). Including the near part gives
\(Lu(x)/u(x)\ge H/512-28\).
At least one of \(l_-,l_+\) is at least one, so
\(H\ge\log(1/(4s))\). This proves \eqref{eq:counter-u9}, with room in its constant.

Every gap has length at most \(|U|\le2\rho\), so \(r(x)\le\rho/2\).
Our choice of \(\rho\) ensures \(T\ge16\), and \eqref{eq:counter-u9} yields
\begin{equation}\label{eq:counter-u14}
\frac{Lu(x)}{u(x)}
 \ge 2(A+40)-\frac{\log2}{512}-32>A+2
 \qquad(x\in U).
\end{equation}
It also shows that \(Lu(x)/u(x)\to+\infty\) as \(x\) approaches either endpoint of any fixed gap.

\subsection{Smooth positive functions on the whole interval}\label{subsec:counter-smooth}

Let \(K_N\) be the union of the closed middle halves of all gaps
removed in the first \(N\) stages. This is a compact subset of \(U\).
For any smooth \(0\le\chi\le1\), vanishing near \(F\) and equal to
one on \(K_N\), put
\[
 A_\chi(x)=\int_U\frac{\chi(y)}{|x-y|}\,dy\qquad(x\in F).
\]
The estimate proving \eqref{eq:counter-u6}, restricted to these middle
halves, gives at least \(1/16\) from each of the first \(N\) stages.
Consequently,
\begin{equation}\label{eq:counter-u15}
 A_\chi(x)\ge N/16\qquad(x\in F).
\end{equation}

Let \(\epsilon_j=1/j\), \(j\ge2\), and put
\[
 d_j=\exp(8-(2j)^{256}),\qquad C_j=2(2j)^{256}.
\]
Since \(r\le D\) and \(p(d_j)=1/(2j)\), one has
\(u(y)\le\epsilon_j/2\) when \(D(y)\le d_j\).
Using this for \(|x-y|<d_j\), and \(u\le1\) on the remainder, gives
\begin{equation}\label{eq:counter-u16}
 \int_U\frac{\chi(y)u(y)}{|x-y|}\,dy
 \le\frac{\epsilon_j}{2}A_\chi(x)+C_j
 \qquad(x\in F).
\end{equation}
Here the far integral is at most \(2\log(2/d_j)\le C_j\).

Set
\[
 k_j=\lceil32(A+1)\rceil+64j(2j)^{256}+j.
\]
Choose a smooth function \(f:\mathbb R\to[0,1]\) with
\(\{f>0\}=U\), and a nondecreasing smooth
\(q_0:\mathbb R\to[0,1]\), equal to zero on \((-\infty,1]\)
and one on \([2,\infty)\). By compactness, choose
\(0<\tau_j\le1/j\) such that \(2\tau_j\le f\) on \(K_{k_j}\),
and define
\[
 \chi_j(x)=q_0(f(x)/\tau_j).
\]
Thus \(\chi_j\) is smooth, vanishes on a neighborhood of \(F\),
and equals one on \(K_{k_j}\). Since \(\tau_j\to0\), at each
fixed point of \(U\) it is eventually identically one on a
neighborhood of that point. Equation~\eqref{eq:counter-u15} gives
\[
 \frac{\epsilon_j}{2}A_{\chi_j}(x)-C_j
 \ge\frac{k_j}{32j}-C_j\ge(A+1)\epsilon_j
 \qquad(x\in F).
\]
Define
\begin{equation}\label{eq:counter-u17}
 v_j=\chi_j u+(1-\chi_j)\epsilon_j.
\end{equation}
This function is strictly positive and smooth on a neighborhood of
\(I\): the function \(u\) is smooth and positive in each gap, and
\(v_j=\epsilon_j\) near \(F\).

At \(x\in F\), \eqref{eq:counter-u16} gives
\(Lv_j(x)\ge(A+1)v_j(x)\). At every fixed \(x\in U\), eventually
\(v_j=u\) on a neighborhood of \(x\). Also \(v_j\to u\) in
\(L^1(I)\), by bounded convergence. The integrals outside that
neighborhood then converge, so
\begin{equation}\label{eq:counter-u18}
 \frac{Lv_j(x)}{v_j(x)}\longrightarrow\frac{Lu(x)}{u(x)}
 >A+2\qquad(x\in U).
\end{equation}
We have therefore constructed positive smooth functions such that
\begin{equation}\label{eq:counter-u19}
 \frac{Lv_j(x)}{v_j(x)}\ge A+1
 \quad\text{eventually for every fixed }x\in I.
\end{equation}
The eventual index is allowed to depend on \(x\).

\subsection{From amplitudes to interpolation nodes}\label{subsec:counter-amplitude}


We prove an upper estimate with an absolute additive constant.

\begin{lemma}\label{lem:amplitude} Let \(h(z)=\sum_{r=0}^{d}h_rz^r\) have real coefficients, no zero on the closed unit disk, and \(h(0)>0\). Set
\[
 v(\cos\theta)=|h(e^{i\theta})|^{-1}.
\]
In this subsection, \(T_k\) denotes the Chebyshev polynomial of the first
kind, defined by \(T_k(\cos\theta)=\cos(k\theta)\).
For \(n>d\), define the real polynomial
\begin{equation}\label{eq:counter-u25}
P_n(x)=\sum_{r=0}^{d}h_rT_{n-r}(x).
\end{equation}
For all sufficiently large \(n\), this polynomial has exactly \(n\) simple roots in \((-1,1)\). The corresponding Lebesgue function satisfies, on any fixed \(K\Subset(-1,1)\),
\begin{equation}\label{eq:counter-u26}
\lambda_n(x)\le\frac2\pi\log n+3-\frac{Lv(x)}{\pi v(x)}
 \qquad(x\in K)
\end{equation}
for all sufficiently large \(n\), with the threshold allowed to depend on \(h,K\).

\end{lemma}

\begin{proof} The analytic logarithm of \(h\), chosen real at zero, is real on the real diameter. Its imaginary part
\(\psi(\theta)=\arg h(e^{i\theta})\) thus satisfies \(\psi(0)=\psi(\pi)=0\). We have the exact identity
\[
 P_n(\cos\theta)=|h(e^{i\theta})|\cos(n\theta-\psi(\theta)).
\]
For \(n>\|\psi'\|_\infty\), the phase increases strictly from \(0\) to \(n\pi\). The polynomial has degree exactly \(n\), since its leading term comes from \(h_0T_n\), with \(h_0>0\). This gives the asserted \(n\) simple roots. At their angles \(\theta_k\),
\begin{equation}\label{eq:counter-u27}
|P_n'(\cos\theta_k)|
 =|h(e^{i\theta_k})|
       \frac{n-\psi'(\theta_k)}{\sin\theta_k}.
\end{equation}

Put \(s=(n\theta-\psi(\theta))/n\), let \(\theta=\Theta_n(s)\) be the inverse, and set \(X_n(s)=\cos\Theta_n(s)\). The roots correspond to
\(s_k=(k-\tfrac12)\pi/n\). Moreover \(X_n\to\cos\) in \(C^3([0,\pi])\). At a point \(x=X_n(s)\) other than a node, put \(b=|\cos(ns)|\). Formula \eqref{eq:counter-u27} gives
\begin{equation}\label{eq:counter-u28}
\lambda_n(x)=\frac{b}{nv(x)}
  \sum_{k=1}^n
  \frac{v(X_n(s_k))|X_n'(s_k)|}{|X_n(s)-X_n(s_k)|}.
\end{equation}

Subtract the simple singularity by defining
\[
 R_n(s,t)=
 \frac{v(X_n(t))|X_n'(t)|}{|X_n(s)-X_n(t)|}
 -\frac{v(X_n(s))}{|s-t|}.
\]
For \(s\) in a fixed compact subinterval of \((0,\pi)\), this function of \(t\) has uniformly bounded variation, with at most a bounded jump at \(t=s\). To verify this near \(s\), factor
\(X_n(t)-X_n(s)=(t-s)D_n(s,t)\). Both \(D_n\) and \(X_n'\) are uniformly bounded away from zero there; the numerator remaining after subtraction vanishes at \(t=s\), so division leaves a uniformly \(C^1\) function on each side, multiplied by \(\operatorname{sgn}(t-s)\). Away from \(s\), the denominator stays uniformly away from zero. Midpoint quadrature therefore has error \(O(1/n)\), uniformly also when \(s\) is arbitrarily close to a quadrature point.

Changing variables \(y=X_n(t)\) in the integral, and first deleting \((s-\epsilon,s+\epsilon)\), gives
\begin{equation}\label{eq:counter-u29}
\begin{split}
 \int_0^\pi R_n(s,t)\,dt
 &=-Lv(x)+v(x)
   \log\frac{1-x^2}{X_n'(s)^2s(\pi-s)}\\
 &=-Lv(x)+v(x)\left[
  2\log(1-\psi'(\theta)/n)-\log(s(\pi-s))\right].
\end{split}
\end{equation}
The cancellation follows because the deleted distances in the \(y\) variable are
\(|X_n'(s)|\epsilon+O(\epsilon^2)\) on both sides.

Write
\[
 \tau=\frac{ns}{\pi}+\frac12=m+t,\qquad m=\lfloor\tau\rfloor,
 \quad0<t<1.
\]
Since \((s-s_k)n/\pi=\tau-k\), midpoint quadrature in \eqref{eq:counter-u28} gives explicitly
\[
 \lambda_n(x)=\frac b\pi\left[
 \sum_{k=1}^n\frac1{|\tau-k|}
 +\frac1{v(x)}\int_0^\pi R_n(s,t')\,dt'\right]+O(1/n).
\]
On the compact interval under consideration, \(1\le m\le n-1\) for large \(n\), and \(b=\sin(\pi t)\). Splitting the elementary harmonic sum at \(m\) gives
\[
 \sum_{k=1}^n\frac1{|\tau-k|}
 \le \log(m(n-m))+2+\frac1t+\frac1{1-t}.
\]
Also, uniformly in this compact interval,
\[
 \log(m(n-m))-\log(s(\pi-s))
 =2\log n-2\log\pi+O(1/n).
\]
Substituting these estimates and \eqref{eq:counter-u29} into \eqref{eq:counter-u28}, and using
\(b/t\le\pi\), \(b/(1-t)\le\pi\), yields
\[
 \lambda_n(x)\le
 b\left[\frac2\pi\log n-\frac{Lv(x)}{\pi v(x)}
                  +\frac{2-2\log\pi}{\pi}\right]+2+O(1/n).
\]
For large \(n\) the bracket is nonnegative, uniformly on \(K\). We may replace \(b\le1\) by one. Since
\(2+(2-2\log\pi)/\pi<3\), this proves \eqref{eq:counter-u26} away from the nodes. At the nodes, \(\lambda_n=1\), and \eqref{eq:counter-u26} follows for large \(n\) as well. \end{proof}

\begin{remark}[A numerical check of the finite term]
At phase midpoints, \(b=1\), and the harmonic sum above satisfies
\[
 \sum_{k=1}^n\frac1{|m+\tfrac12-k|}
 =\log(m(n-m))+2\gamma+4\log2+o(1),
\]
uniformly on compact interior intervals. Thus the corrected quantity
\[
 E_n(x)=\lambda_n(x)-\frac2\pi\log n+\frac{Lv(x)}{\pi v(x)}
\]
at these points tends to
\(\kappa_V=(2/\pi)(\gamma+\log(4/\pi))=0.521251626\ldots\),
the constant in V\'ertesi's asymptotic formula \cite{Vertesi1990}.
The reproducible script \texttt{checks/amplitude\_numerics.py} evaluates
three fixed examples \(h=1,\ 1-z/2,\ (1-z/2)^2\) for
\(n=50,100,\ldots,800\), using the derivative-weight identity and
closed expressions for \(Lv/v\). On the prescribed phase midpoints
in \(K=[-0.8,0.8]\), the largest computed values at \(n=800\) are,
respectively, \(0.5212524\), \(0.5215144\), and \(0.5217777\).
\end{remark}

Every positive smooth function \(v\) on \(I\) can be approximated in
\(C^1(I)\) by amplitudes occurring in this lemma. Approximate \(1/v\)
by a strictly positive real polynomial \(q\) of degree \(d\), and form
\[
 J(z)=z^d q\bigl((z+z^{-1})/2\bigr).
\]
This is a real polynomial, and \(|J(e^{i\theta})|=q(\cos\theta)>0\).
In its factorization, replace every factor \(z-r\) with \(|r|<1\)
by \(1-\overline r z\), retaining the leading scalar. The resulting
polynomial \(h\) has no zero on the closed disk and the same modulus
as \(J\) on the unit circle. Conjugate roots are reflected together,
so its coefficients are real; changing its sign makes \(h(0)>0\).
Thus
\[
 |h(e^{i\theta})|=q(\cos\theta),\qquad
 |h(e^{i\theta})|^{-1}=q(\cos\theta)^{-1}.
\]
Polynomial approximation in \(C^1\) follows by uniformly approximating
the derivative and then integrating. Inversion is continuous in \(C^1\)
on functions with a positive lower bound. Hence \(q^{-1}\) approximates
\(v\) in \(C^1(I)\), and \eqref{eq:counter-u3} gives arbitrarily accurate
uniform approximation of \(Lv/v\).


\begin{remark}[Dependence on interval length]
The loss on a small interval is already visible in a simpler construction.
Let \(I_\ell=[-\ell/2,\ell/2]\), where \(0<\ell<1/4\), and choose a
smooth \(v\) equal to \(1\) on \(I_\ell\) and to
\(\delta=1/\log(1/\ell)\) outside \([-\ell,\ell]\), with
\(\delta\le v\le1\) and \(\|v'\|_\infty=O(1/\ell)\).
The integral outside \([-\ell,\ell]\) and the transition regions give,
uniformly on \(I_\ell\),
\[
 Lv/v=2\log(1/\ell)+O(1).
\]
Approximate by an amplitude as above, with error at most one in this
quotient. Lemma~\ref{lem:amplitude} gives the upper bound in
\[
 \sup_{I_\ell}\lambda_n=\frac2\pi\log(n\ell)+O(1).
\]
The matching lower bound follows from \eqref{eq:counter-u28}--\eqref{eq:counter-u29}
at a phase midpoint, where \(b=1\); such points lie in \(I_\ell\) for
all large \(n\). The constants here can be absolute, with the starting
index depending on \(\ell\).
Thus Tao's interval constant can be at least
\((2/\pi)\log(1/|I_\ell|)-O(1)\) when the array changes with \(\ell\).
The gap losses in Theorem~\ref{thm:nonuniform} are consistent with this
dependence; its additional conclusion concerns one fixed array for
which every finite common constant fails.
\end{remark}

\subsection{Assembly of the array}\label{subsec:counter-array}

\begin{proof}[Proof of Theorem~\ref{thm:nonuniform}]

Fix \(M>0\), and carry out Subsections~\ref{subsec:counter-set}--\ref{subsec:counter-smooth} with \(A=\pi(M+4)\).
For each \(j\ge2\), choose an amplitude \(\widetilde v_j=|h_j|^{-1}\) as in Subsection~\ref{subsec:counter-amplitude} so close to \(v_j\) that
\begin{equation}\label{eq:counter-u30}
\left\|\frac{L\widetilde v_j}{\widetilde v_j}
                 -\frac{Lv_j}{v_j}\right\|_\infty<1/2.
\end{equation}
Let \(K_j=[-1+1/j,1-1/j]\). Choose a strictly increasing sequence of integers \(N_j\), each larger than \(\deg h_j\) and the threshold in Lemma~\ref{lem:amplitude} for \(h_j,K_j\). For \(N_j\le n<N_{j+1}\), use as row \(n\) the \(n\) roots of \eqref{eq:counter-u25} with \(h=h_j\). Fill the finitely many earlier rows with Chebyshev roots.

Fix \(x\in(-1,1)\). The block index \(j\) tends to infinity with \(n\), and \(x\in K_j\) for all large \(j\). Equations \eqref{eq:counter-u19} and \eqref{eq:counter-u30} give
\(L\widetilde v_j(x)/\widetilde v_j(x)\ge A+1/2\) eventually. Hence \eqref{eq:counter-u26} gives, for all sufficiently large \(n\),
\[
 \lambda_n(x)\le\frac2\pi\log n+3-\frac{A+1/2}{\pi}
 <\frac2\pi\log n-M.
\]
This proves \eqref{eq:counter-eventual} and the failure of an absolute additive constant.

For the final assertion, \eqref{eq:counter-u18}, \eqref{eq:counter-u26}, and \eqref{eq:counter-u30} give at every fixed \(x\in U\)
\begin{equation}\label{eq:counter-u31}
\limsup_{n\to\infty}
 \left(\lambda_n(x)-\frac2\pi\log n\right)
 \le3+\frac1{2\pi}-\frac{Lu(x)}{\pi u(x)}.
\end{equation}
Choose a gap \((b,c)\) with both endpoints in \((-1,1)\). By \eqref{eq:counter-u9}, \(Lu(x)/u(x)\to+\infty\) as \(x\downarrow b\). Given any finite \(C\), there is a nonempty interval \((b,b+\sigma)\) on which the right side of \eqref{eq:counter-u31} is strictly less than \(-C-1\). This interval is disjoint from \(G_C\), proving the second assertion. \end{proof}


\section{Consequences}
\label{sec:consequences}

\begin{corollary}[Positive-measure sets]\label{cor:positive-measure}
Let \(E\subset(-1,1)\) be measurable with \(|E|>0\), and let
\(0<c<2/\pi\). For all sufficiently large \(n\), there is \(x\in E\)
such that \(\lambda_n(x)>c\log n\).
\end{corollary}
\begin{proof}
Otherwise \(\lambda_n(y)\le c\log n\) for every \(y\in E\) and every
\(n\) in an infinite set \(S\). The argument in Section~\ref{sec:ae} applies to these
rows, with all limits taken as \(n\to\infty\) through \(S\), and gives
the same contradiction \(|E|/A\ge|E|\) with \(A>1\).
\end{proof}

\begin{corollary}[No uniform interval constant]\label{cor:no-uniform-interval}
For the array satisfying both assertions of Theorem~\ref{thm:nonuniform},
with any \(M>0\), there is no finite \(C\) such that
\[
 \sup_{x\in I}\lambda_n(x)\ge\frac2\pi\log n-C
\]
for every nondegenerate compact interval \(I\subset(-1,1)\) and all sufficiently
large \(n\), with the threshold allowed to depend on \(I\).
\end{corollary}
\begin{proof}
Suppose that such a \(C\) exists. If \(G_{C+1}\) misses a nondegenerate
compact interval \(K_0\subset(-1,1)\), the closed sets
\[
 F_N=\{x\in K_0:\lambda_n(x)\le(2/\pi)\log n-C-1
                    \text{ for all }n\ge N\}
\]
cover \(K_0\). Baire's theorem gives an \(F_N\) containing a
nondegenerate compact subinterval \(I\). The resulting upper bound on \(I\)
contradicts the assumed lower bound for large \(n\).
Thus \(G_{C+1}\) is dense, contrary to Theorem~\ref{thm:nonuniform}.
\end{proof}

\section*{Use of generative AI}

GPT-6 Astra was used to generate proofs and draft the manuscript;
GPT-5.6 Sol and Claude Opus 5 were used for editorial review.
The Lean formalization was generated using OpenAI Codex (GPT-6).

\begin{thebibliography}{11}

\bibitem{Bernstein}
S.~Bernstein,
Sur la limitation des valeurs d'un polyn\^ome \(P_n(x)\) de degr\'e \(n\)
sur tout un segment par ses valeurs en \((n+1)\) points du segment,
\emph{Izv. Akad. Nauk SSSR}, Ser.~VII (1931), no.~8, 1025--1050.

\bibitem{Bloom}
T.~F. Bloom,
\emph{Erd\H{o}s Problem \#1132},
\url{https://www.erdosproblems.com/1132}.
Accessed 6 September 2026.

\bibitem{Brutman}
L.~Brutman,
Lebesgue functions for polynomial interpolation---a survey,
\emph{Ann. Numer. Math.} \textbf{4} (1997), 111--127.

\bibitem{Erdos}
P.~Erd\H{o}s,
Problems and results on the convergence and divergence properties of
the Lagrange interpolation polynomials and some extremal problems,
\emph{Mathematica (Cluj)} \textbf{10} (1968), 65--73.
\href{https://combinatorica.hu/~p_erdos/1967-20.pdf}{Original paper}.

\bibitem{Erdos1961}
P.~Erd\H{o}s,
Problems and results on the theory of interpolation. II,
\emph{Acta Math. Acad. Sci. Hungar.} \textbf{12} (1961), 235--244.

\bibitem{ErdosTuran}
P.~Erd\H{o}s and P.~Tur\'an,
An extremal problem in the theory of interpolation,
\emph{Acta Math. Acad. Sci. Hungar.} \textbf{12} (1961), 221--234.

\bibitem{ErdosVertesi}
P.~Erd\H{o}s and P.~V\'ertesi,
On the almost everywhere divergence of Lagrange interpolatory polynomials
for arbitrary system of nodes,
\emph{Acta Math. Acad. Sci. Hungar.} \textbf{36} (1980), 71--89.

\bibitem{ErdosVertesi1981}
P.~Erd\H{o}s and P.~V\'ertesi,
On the Lebesgue function of interpolation,
in \emph{Functional Analysis and Approximation} (Oberwolfach, 1980),
P.~L.~Butzer, B.~Sz.-Nagy and E.~G\"orlich (eds.),
International Series of Numerical Mathematics \textbf{60},
Birkh\"auser, Basel, 1981, 299--309.
\href{https://doi.org/10.1007/978-3-0348-9369-5_27}{doi:10.1007/978-3-0348-9369-5\_27}.

\bibitem{Faber}
G.~Faber,
\"Uber die interpolatorische Darstellung stetiger Funktionen,
\emph{Jahresber. Deutsch. Math.-Verein.} \textbf{23} (1914), 192--210.

\bibitem{Tao}
T.~Tao,
\emph{Local Bernstein theory, and lower bounds for Lebesgue constants},
\href{https://arxiv.org/abs/2603.21453v3}{arXiv:2603.21453v3}
(2026).

\bibitem{Vertesi1990}
P.~V\'ertesi,
Optimal Lebesgue constant for Lagrange interpolation,
\emph{SIAM J. Numer. Anal.} \textbf{27} (1990), 1322--1331.
\href{https://doi.org/10.1137/0727075}{doi:10.1137/0727075}.

\end{thebibliography}

\end{document}
